\section*{Capítulo \ref{ch:g11:fossil-fuel-geology} --- A geologia do carvão, do petróleo e do gás}

\begin{solution}{exo:g11:fossil-fuel-geology:1}
O plâncton morto se deposita num fundo de mar pobre em oxigênio; a matéria
orgânica vira \omterm{def:g11:fossil-fuel-geology:source-rock}{querogênio} numa lama; o soterramento sob sedimentos mais
jovens; a \omterm{def:g11:fossil-fuel-geology:oil-window}{maturação} na \omterm{def:g11:fossil-fuel-geology:oil-window}{janela do petróleo}; a \omterm{def:g11:fossil-fuel-geology:migration}{migração}; o petróleo se acumula
sob uma \omterm{def:g11:fossil-fuel-geology:reservoir}{rocha selante}.
\end{solution}

\begin{solution}{exo:g11:fossil-fuel-geology:2}
(1) \omterm{def:g11:fossil-fuel-geology:source-rock}{Rocha geradora}. (2) \omterm{def:g11:fossil-fuel-geology:reservoir}{Rocha reservatório}. (3) \omterm{def:g11:fossil-fuel-geology:reservoir}{Rocha selante}. (4) \omterm{def:g11:fossil-fuel-geology:reservoir}{Rocha
selante}. (5) \omterm{def:g11:fossil-fuel-geology:reservoir}{Rocha reservatório}.
\end{solution}

\begin{solution}{exo:g11:fossil-fuel-geology:3}
Topo: $(60 - 10)/25 = \qty{2.0}{km}$; base: $(120 - 10)/25 = \qty{4.4}{km}$.
A janela tem $60/25 = \qty{2.4}{km}$ de espessura.
\end{solution}

\begin{solution}{exo:g11:fossil-fuel-geology:4}
(1) Turfa, linhito, \omterm{def:g11:fossil-fuel-geology:coal}{carvão} betuminoso, antracito. (2) O antracito, de 86 a
\qty{97}{\percent} de carbono. (3) O antracito, além de cerca de
\qty{5}{km} e de \qty{150}{\degreeCelsius}.
\end{solution}

\begin{solution}{exo:g11:fossil-fuel-geology:5}
(1) Os fluidos se separam pela densidade: o gás é o mais leve, depois o
petróleo, depois a água. (2) O petróleo e o gás sobem porque são mais leves
que a água; só uma rocha que eles não conseguem atravessar os detém. O
lamito (ou o folhelho) e o sal.
\end{solution}

\begin{solution}{exo:g11:fossil-fuel-geology:6}
(1) A: $(60 - 15)/20 = \qty{2.25}{km}$ e $(120 - 15)/20 = \qty{5.25}{km}$.
B: $45/40 \approx \qty{1.13}{km}$ e $105/40 \approx \qty{2.63}{km}$.
(2) A \qty{2}{km}: em A, $15 + 40 = \qty{55}{\degreeCelsius}$, ainda
imatura, nada de petróleo; em B, $15 + 80 = \qty{95}{\degreeCelsius}$, na
\omterm{def:g11:fossil-fuel-geology:oil-window}{janela do petróleo}, ela produz petróleo. (3) Na bacia A: ali o petróleo só
se forma abaixo de \qty{2.25}{km}.
\end{solution}

\begin{solution}{exo:g11:fossil-fuel-geology:7}
(1) $11 \times 3 = \qty{33}{m}$ de turfa. (2) $\qty{33000}{mm}$ a
\qty{1}{mm} por ano: cerca de \num{33000} anos. (3) A turfa só se conserva
onde fica sob água parada, fora do alcance do oxigênio; à medida que ela se
acumula, o terreno precisa afundar no mesmo ritmo para que a sua superfície
continue encharcada.
\end{solution}

\begin{solution}{exo:g11:fossil-fuel-geology:8}
(1) $20 + 30 \times 1.2 = \qty{56}{\degreeCelsius}$;
$20 + 30 \times 2.3 = \qty{89}{\degreeCelsius}$;
$20 + 30 \times 4.0 = \qty{140}{\degreeCelsius}$.
(2) Há 120 milhões de anos, imatura; há 95 milhões de anos, na \omterm{def:g11:fossil-fuel-geology:oil-window}{janela do
petróleo}; há 20 milhões de anos, na \omterm{def:g11:fossil-fuel-geology:oil-window}{janela do gás}.
\end{solution}

\begin{solution}{exo:g11:fossil-fuel-geology:9}
(1) $2580 - 2400 = \qty{180}{m}$. (2) O gás ocupa
$2450 - 2400 = \qty{50}{m}$, o petróleo $2580 - 2450 = \qty{130}{m}$. (3) Eles
transbordam no ponto de derrame, abaixo da curva de nível fechada mais
profunda, e continuam a migrar para cima, até outra \omterm{def:g11:fossil-fuel-geology:trap}{armadilha} ou até a
superfície.
\end{solution}

\begin{solution}{exo:g11:fossil-fuel-geology:10}
(1) Topo da janela a \qty{3}{km}: $(60 - 20)/3 \approx
\qty{13}{\degreeCelsius/km}$; base: $(120 - 20)/3 \approx
\qty{33}{\degreeCelsius/km}$. (2) A \qty{3}{km}, a temperatura vai de
$20 + 75 = \qty{95}{\degreeCelsius}$ a $20 + 102 = \qty{122}{\degreeCelsius}$:
na \omterm{def:g11:fossil-fuel-geology:oil-window}{janela do petróleo}, ou na sua base para os gradientes mais íngremes. A
regra é razoável: na maioria das bacias, uma \omterm{def:g11:fossil-fuel-geology:source-rock}{rocha geradora} sob \qty{3}{km}
de rochas já começou a produzir petróleo.
\end{solution}

\begin{solution}{exo:g11:fossil-fuel-geology:11}
(1) \qty{30}{\percent}: linhito; \qty{70}{\percent}: \omterm{def:g11:fossil-fuel-geology:coal}{carvão} betuminoso;
\qty{92}{\percent}: antracito. (2) O antracito, além de cerca de
\qty{5}{km} e de \qty{150}{\degreeCelsius}. (3) Preto e vítreo.
\end{solution}

\begin{solution}{exo:g11:fossil-fuel-geology:12}
(1) $1100/25 = \qty{44}{m}$ por milhão de anos. (2) $0.044 \times 30
\approx \qty{1.3}{\degreeCelsius}$ por milhão de anos. (3) A janela cobre
\qty{60}{\degreeCelsius}: $60/1.32 \approx 45$ milhões de anos. (4) O
soterramento ficou mais lento depois de 95 milhões de anos atrás (cerca de
24, depois 23 e \qty{20}{m} por milhão de anos): a \omterm{def:g11:fossil-fuel-geology:source-rock}{rocha geradora} se
aqueceu mais devagar.
\end{solution}

\begin{solution}{exo:g11:fossil-fuel-geology:13}
(1) $(80 - 51)/(117 - 51) = 29/66 \approx 0.44$: cerca de \qty{44}{\percent}.
(2) Ele migrou para cima antes que existisse qualquer \omterm{def:g11:fossil-fuel-geology:trap}{armadilha} acima: escapou
para a superfície, onde foi oxidado, ou se acumulou em alguma outra
\omterm{def:g11:fossil-fuel-geology:trap}{armadilha}. (3) Há 40 milhões de anos, a \omterm{def:g11:fossil-fuel-geology:source-rock}{rocha geradora} já tinha saído da
\omterm{def:g11:fossil-fuel-geology:oil-window}{janela do petróleo}: ela só produzia gás, e o anticlinal só poderia conter
gás vindo dela.
\end{solution}

\begin{solution}{exo:g11:fossil-fuel-geology:14}
(1) Topo: $(60 - 20)/40 = \qty{1.0}{km}$; base: $(120 - 20)/40 =
\qty{2.5}{km}$. (2) \qty{1.0}{km} fica entre \qty{0.4}{km} (há 150 milhões
de anos) e \qty{1.2}{km} (120): $150 - 30 \times 0.6/0.8 = 127.5$
milhões de anos atrás. \qty{2.5}{km} fica entre \qty{2.3}{km} (95) e
\qty{3.0}{km} (66): $95 - 29 \times 0.2/0.7 \approx 86.7$ milhões de anos
atrás. (3) Nada de petróleo: a \omterm{def:g11:fossil-fuel-geology:source-rock}{rocha geradora} saiu da \omterm{def:g11:fossil-fuel-geology:oil-window}{janela do petróleo}
antes que o anticlinal se formasse, e o seu petróleo já tinha migrado para
longe. Ele poderia conter o gás produzido depois de 80 milhões de anos atrás.
\end{solution}

\begin{solution}{exo:g11:fossil-fuel-geology:15}
(1) Pelo menos \qty{5}{km}. (2) $15 + 30 \times 5 = \qty{165}{\degreeCelsius}$,
acima de \qty{150}{\degreeCelsius}: coerente. (3) Pelo menos \qty{5}{km} de
rochas foram soerguidos e erodidos. $\qty{5e6}{mm}$ em no máximo
\num{299} milhões de anos dá pelo menos $5 \times 10^6/2.99 \times 10^8
\approx \qty{0.017}{mm}$ por ano em média.
\end{solution}

\begin{solution}{pb:g11:fossil-fuel-geology:1}
\textbf{1.} \omterm{def:g11:fossil-fuel-geology:source-rock}{Rocha geradora}: o folhelho negro; reservatório: o arenito poroso;
\omterm{def:g11:fossil-fuel-geology:reservoir}{rocha selante}: o lamito acima dele.

\textbf{2.} O plâncton que se depositou num fundo de mar pobre em oxigênio
escapou da decomposição e foi enterrado na lama.

\textbf{3.} A \omterm{def:g11:fossil-fuel-geology:source-rock}{rocha geradora} no Jurássico (165 fica entre 201.4 e
143.1). O arenito começou no Jurássico e foi depositado principalmente no
Cretáceo Inferior.

\textbf{4.} A matéria orgânica dispersa numa rocha sedimentar, insolúvel em
solventes orgânicos. Ela só produz petróleo depois de aquecida a cerca de
\qty{60}{\degreeCelsius}, o que exige soterramento.

\textbf{5.} $400 + 800 + 1100 + 700 + 700 + 300 - 500 = \qty{3500}{m}$.

\textbf{6.} \qty{0.4}{km}, \qty{1.2}{km}, \qty{2.3}{km}, \qty{3.0}{km},
\qty{3.7}{km}, \qty{4.0}{km}; hoje \qty{3.5}{km}.

\textbf{7.} $400/15 \approx 26.7$; $800/30 \approx 26.7$; $1100/25 = 44$;
$700/29 \approx 24.1$; $700/31 \approx 22.6$; $300/15 = 20$ (\unit{m} por
milhão de anos).

\textbf{8.} Há 20 milhões de anos, a \qty{4.0}{km}.

\textbf{9.} A bacia foi soerguida e \qty{500}{m} de rochas foram erodidos do
seu topo.

\textbf{10.} A base do arenito foi depositada há 150 milhões de anos, sobre
\qty{400}{m} de folhelho e de lamito: $3.5 - 0.4 = \qty{3.1}{km}$.

\textbf{11.} $T = 20 + 30 z$: \qty{32}{\degreeCelsius}, \num{56}, \num{89},
\num{110}, \num{131}, \qty{140}{\degreeCelsius}; hoje
\qty{125}{\degreeCelsius}.

\textbf{12.} $40/30 \approx \qty{1.33}{km}$ e $100/30 \approx
\qty{3.33}{km}$.

\textbf{13.} Entre há 120 milhões de anos (\qty{56}{\degreeCelsius}) e há
95 milhões de anos (\qty{89}{\degreeCelsius}).

\textbf{14.} A temperatura subiu 33 graus em 25 milhões de anos; ela levou
$25 \times 4/33 \approx 3.0$ milhões de anos para subir os 4 graus de
\num{56} a \qty{60}{\degreeCelsius}: há cerca de $120 - 3.0 = 117$ milhões
de anos.

\textbf{15.} Entre há 66 milhões de anos (\qty{110}{\degreeCelsius}) e há
35 (\qty{131}{\degreeCelsius}): $66 - 31 \times 10/21 \approx 51.2$ milhões
de anos atrás.

\textbf{16.} \qty{140}{\degreeCelsius}, há 20 milhões de anos. A
\qty{125}{\degreeCelsius} hoje, ela está na \omterm{def:g11:fossil-fuel-geology:oil-window}{janela do gás}: produz
principalmente metano.

\textbf{17.} Sem nenhuma \omterm{def:g11:fossil-fuel-geology:trap}{armadilha} acima dele, ele migrou para cima e
escapou, para a superfície ou para outras \omterm{def:g11:fossil-fuel-geology:trap}{armadilhas}.

\textbf{18.} $(80 - 51.2)/(117.0 - 51.2) \approx 0.44$.

\textbf{19.} Gás em cima, depois petróleo, depois água.

\textbf{20.} A \omterm{def:g11:fossil-fuel-geology:source-rock}{rocha geradora} entrou na \omterm{def:g11:fossil-fuel-geology:oil-window}{janela do petróleo} há cerca de
\textbf{117 milhões de anos}.
\end{solution}
